\documentclass{article} % For LaTeX2e
\usepackage{iclr2026_conference,times}

% Optional math commands from https://github.com/goodfeli/dlbook_notation.
\input{math_commands.tex}

\usepackage{hyperref}
\usepackage{url}

\title{Hardware-Aware Expert Placement for Reducing Intra-GPU Communication in Mixture-of-Experts Models}

% Keep commented for anonymous submission.
\iclrfinalcopy

\author{Jiahao Chen \\
Georgia Institute of Technology \\
\texttt{jchen3392@gatech.edu}
\And
Jane Yijun Sun \\
Georgia Institute of Technology \\
\texttt{ysun713@gatech.edu}
}

\begin{document}

\maketitle

\begin{abstract}
This project extends collaborative communication ideas for mixture-of-experts (MoE) systems from inter-GPU communication to intra-GPU optimization. The core question is whether hardware-aware expert placement can reduce memory traffic and execution latency within a single GPU. Building on Occult (Luo et al., ICML 2025), we propose to model expert co-activation patterns, design placement strategies that keep frequently co-activated experts close in the execution pipeline, and evaluate performance using latency, throughput, utilization, and profiling-based memory metrics. We additionally provide an updated implementation plan based on DeepSeek-MoE for practical evaluation.
\end{abstract}

\section{Research Paper Summary}
This project builds upon the paper ``Occult: Optimizing Collaborative Communication across Experts for Accelerated Parallel MoE Training and Inference'' (Luo et al., ICML 2025). Mixture-of-Experts (MoE) models scale large language models by activating only a subset of experts per token. This creates communication overhead because tokens are routed to different experts, often across GPUs. Occult addresses this by modeling expert co-activation and placing frequently collaborating experts on the same device to reduce cross-device communication.

\section{Summary of Contributions of the Paper}
Occult introduces collaborative communication to estimate expert co-activation patterns and optimize expert placement across devices. The method increases intra-device collaboration and reduces token duplication. The paper also proposes collaboration pruning to limit routing choices and custom sparse matrix multiplication kernels to improve MoE efficiency. Reported results show substantial speedups for distributed MoE training and inference.

\section{Limitations}
Occult mainly optimizes inter-GPU communication and treats intra-GPU communication cost as roughly uniform. This overlooks architectural hierarchy inside modern GPUs, including streaming multiprocessors (SMs), shared memory, and cache levels. As a result, token movement inside a single GPU can still cause non-trivial memory traffic and latency.

\section{Project Goal}
This project investigates whether collaborative communication can be extended to intra-GPU optimization. The central research question is: \emph{Can hardware-aware expert placement reduce memory traffic and execution latency within a single GPU for MoE models?}

We will analyze expert co-activation patterns and design placement strategies that reduce token movement across GPU execution units.

\section{High-Level Methods}
We will first implement a baseline MoE model in PyTorch and collect routing statistics. These statistics will be used to construct an expert collaboration graph that captures how often expert pairs are activated together.

Next, we will implement expert placement policies that place frequently co-activated experts near each other in the GPU execution pipeline. We will compare multiple policies, including random placement and collaboration-aware placement.

\section{High-Level Evaluation}
Evaluation will use quantitative metrics including inference latency, throughput (tokens per second), GPU utilization, and memory traffic. NVIDIA Nsight profiling will be used to measure memory bandwidth use and kernel-level behavior. Results will be compared against baseline implementations without hardware-aware expert placement.

\section{Implementation Plan}
\subsection{Model}
Our experiments will be conducted using the DeepSeek-MoE architecture, an open-source Mixture-of-Experts (MoE) large language model. DeepSeek-MoE replaces the standard feed-forward network in transformer layers with a set of expert networks and a routing module. For each token, the router selects top-$k$ experts, and their outputs are combined to produce the final representation.

We initially plan to run the original DeepSeek-MoE 16B model. This model contains approximately 16 billion parameters and multiple MoE layers with dozens of experts per layer. Using the original model helps capture realistic expert collaboration patterns and token routing behavior in large-scale MoE systems.

Experiments will be implemented in PyTorch and executed on GPUs available through the PACE computing cluster. Our primary focus is the expert execution stage of the MoE pipeline, where tokens are dispatched to selected experts and processed by expert MLP networks. This stage is the dominant source of memory movement and communication overhead.

If the full DeepSeek-MoE 16B model cannot run efficiently within available GPU memory on PACE, we will use a scaled-down configuration that preserves MoE routing and expert structure while reducing computational cost. A representative scaled setup is:
\begin{itemize}
\item 12 transformer layers
\item 16 experts per MoE layer
\item top-2 routing
\item hidden dimension between 1024 and 2048
\end{itemize}

Because this project focuses on system-level behavior rather than model accuracy, retraining is not strictly required. Even with randomly initialized weights or partially loaded pretrained weights, the architecture still produces routing decisions and expert activations that generate realistic token dispatch patterns for studying memory traffic and execution efficiency.

\subsection{Method 1: Collaboration-Aware Expert Placement}
The first method analyzes expert co-activation patterns to identify experts that are frequently selected together by the routing network. During execution, we will collect routing statistics to construct an expert collaboration matrix that records how often pairs of experts are activated together.

Using this matrix, we will design an expert placement strategy that groups frequently co-activated experts within the same execution region. Although explicit control over GPU streaming multiprocessor placement is limited, we can approximate locality by organizing experts into execution groups that run within the same kernel execution sequence.

The intuition is that frequently collaborating experts executed close together may allow tokens to remain in fast GPU memory structures such as caches or shared memory rather than being repeatedly loaded from global memory. We will compare this collaboration-aware strategy against a baseline where experts are executed in random or default order.

\subsection{Method 2: Thread-Block Expert Fusion}
The second method aims to reduce redundant memory access by fusing frequently co-activated experts into a single thread-block execution. In standard MoE implementations, tokens routed to multiple experts may be loaded from global memory multiple times, increasing memory bandwidth usage.

Using the collaboration matrix, we will identify expert pairs that frequently activate together and combine them into fused expert kernels. In this design, tokens are loaded once into registers or shared memory and reused across multiple expert computations before outputs are written back to global memory.

This approach reduces redundant memory loads and may also reduce kernel launch overhead, potentially improving runtime performance.

\subsection{Method 3: Shared-Memory Token Reuse}
The third method focuses on exploiting the GPU memory hierarchy by storing tokens in shared memory during expert computation. Instead of loading tokens from global memory for each expert execution, tokens assigned to a group of experts are first loaded into shared-memory buffers, and experts operate sequentially on this shared representation.

Since shared memory has significantly lower latency and higher bandwidth than global memory, this approach can reduce global memory traffic and improve execution efficiency. Due to limited shared-memory capacity, tokens may be processed in tiles that fit within available space. We will experiment with different tile sizes to determine effective configurations.

\subsection{Implementation Evaluation}
We will evaluate these methods against a baseline MoE implementation without hardware-aware optimizations. Performance metrics include:
\begin{itemize}
\item inference latency
\item throughput (tokens per second)
\item global memory bandwidth usage
\item GPU utilization
\end{itemize}

We will use PyTorch Profiler and NVIDIA Nsight Compute to analyze kernel execution behavior, memory traffic, and GPU resource utilization. These measurements will determine whether hardware-aware expert placement effectively reduces intra-GPU communication overhead and improves execution efficiency.

\nocite{*}
\bibliographystyle{iclr2026_conference}
\bibliography{iclr2026_conference}

\end{document}
